
\documentclass[UTF8,12pt,a4paper]{ctexart}
\usepackage{amsmath,amssymb}
\usepackage{geometry}
\geometry{margin=2cm}

% ===== 水印：贺昱琪学习资料 =====
\usepackage{xcolor}
\usepackage{eso-pic}
\usepackage{graphicx}

\newcommand{\WMText}{贺昱琪学习资料}
\newcommand{\WMScale}{2.28}
\newcommand{\WMAngle}{45}
\colorlet{WMColor}{black!8}

\AddToShipoutPictureBG{%
  \AtPageLowerLeft{%
    \put(\LenToUnit{0.66\paperwidth},\LenToUnit{0.70\paperheight}){%
      \rotatebox{\WMAngle}{\scalebox{\WMScale}{\textcolor{WMColor}{\WMText}}}%
    }%
    \put(\LenToUnit{0.33\paperwidth},\LenToUnit{0.50\paperheight}){%
      \rotatebox{\WMAngle}{\scalebox{\WMScale}{\textcolor{WMColor}{\WMText}}}%
    }%
    \put(\LenToUnit{0.66\paperwidth},\LenToUnit{0.30\paperheight}){%
      \rotatebox{\WMAngle}{\scalebox{\WMScale}{\textcolor{WMColor}{\WMText}}}%
    }%
  }%
}

% ===== 4题铺满一页布局 =====
\usepackage{enumitem}
\newcommand{\BottomWriteSpace}{2.5cm}

\newenvironment{OnePageFourFill}{
  \small
  \vspace*{0pt}
  \begin{enumerate}[leftmargin=*, itemsep=0pt, topsep=0pt, parsep=0pt, partopsep=0pt]
  \setlength{\parskip}{0pt}
}{
  \end{enumerate}
  \vspace*{\BottomWriteSpace}
  \normalsize
  \newpage
}

\newcommand{\Q}[1]{\item #1\par\vfill}
\newcommand{\Qlast}[1]{\item #1\par}

\begin{document}

% =========================================================
% 6.29
% =========================================================
\section*{6月29日作业（4道题当晚22:00前打卡）}
\begin{OnePageFourFill}

\Q{计算：$\displaystyle 2021 \div 2021\frac{2021}{2022} \div \frac{2022 + \dfrac{1}{2023}}{2023 + \dfrac{1}{2022}}$}

\Q{计算：$\displaystyle \frac{5}{7} \times 49\frac{12}{13} + 30 \times \frac{1}{7} + \frac{1}{7} \times \frac{5}{13} + 2\frac{2}{15} \div \frac{4}{45}$}

\Q{（工程问题）一项工程，甲工程队做4天休息1天，乙工程队做5天休息2天。甲工程队单独做完要37天，乙工程队单独做完要54天。现在两队合作，那么完成这项工程要\underline{\hspace{3em}}天。}

\Qlast{（容斥原理）某校 150 名学生在一次语、数、外三科竞赛中，参加语文竞赛的有 70 人，参加数学竞赛的有 75 人，参加外语竞赛的有 65 人，既参加语文又参加数学竞赛的有 25 人，既参加数学又参加外语竞赛的有 30 人，既参加语文又参加外语竞赛的有 20 人，有 5 人这三项竞赛都不参加。则三项都参加的共有多少人？}

\end{OnePageFourFill}

% =========================================================
% 6.30
% =========================================================
\section*{6月30日作业（4道题当晚22:00前打卡）}
\begin{OnePageFourFill}

\Q{计算：$\displaystyle 2022 \div 2022\frac{2022}{2023} \div \frac{2023 + \dfrac{1}{2024}}{2024 + \dfrac{1}{2023}}$}

\Q{计算：$\displaystyle \frac{3}{8} \times 59\frac{10}{11} + 20 \times \frac{1}{8} + \frac{1}{8} \times \frac{3}{11} + 4\frac{1}{6} \div \frac{5}{36}$}

\Q{（工程问题）一项工程，甲工程队做3天休息1天，乙工程队做4天休息1天。甲工程队单独做完要19天，乙工程队单独做完要29天。现在两队合作，那么完成这项工程要\underline{\hspace{3em}}天。}

\Qlast{（容斥原理）某校 220 名学生在一次语、数、外三科竞赛中，参加语文竞赛的有 100 人，参加数学竞赛的有 110 人，参加外语竞赛的有 90 人，既参加语文又参加数学竞赛的有 35 人，既参加数学又参加外语竞赛的有 40 人，既参加语文又参加外语竞赛的有 30 人，有 10 人这三项竞赛都不参加。则三项都参加的共有多少人？}

\end{OnePageFourFill}

% =========================================================
% 7.1
% =========================================================
\section*{7月1日作业（4道题当晚22:00前打卡）}
\begin{OnePageFourFill}

\Q{计算：$\displaystyle 2023 \div 2023\frac{2023}{2024} \div \frac{2024 + \dfrac{1}{2025}}{2025 + \dfrac{1}{2024}}$}

\Q{计算：$\displaystyle \frac{2}{9} \times 89\frac{16}{17} + 90 \times \frac{1}{9} + \frac{1}{9} \times \frac{2}{17} + 3\frac{1}{7} \div \frac{11}{49}$}

\Q{（工程问题）一项工程，甲工程队做5天休息2天，乙工程队做3天休息1天。甲工程队单独做完要26天，乙工程队单独做完要27天。现在两队合作，那么完成这项工程要\underline{\hspace{3em}}天。}

\Qlast{（容斥原理）某校 180 名学生在一次语、数、外三科竞赛中，参加语文竞赛的有 80 人，参加数学竞赛的有 90 人，参加外语竞赛的有 74 人，既参加语文又参加数学竞赛的有 28 人，既参加数学又参加外语竞赛的有 32 人，既参加语文又参加外语竞赛的有 24 人，有 8 人这三项竞赛都不参加。则三项都参加的共有多少人？}

\end{OnePageFourFill}

% =========================================================
% 7.2
% =========================================================
\section*{7月2日作业（4道题当晚22:00前打卡）}
\begin{OnePageFourFill}

\Q{计算：$\displaystyle 2024 \div 2024\frac{2024}{2025} \div \frac{2025 + \dfrac{1}{2026}}{2026 + \dfrac{1}{2025}}$}

\Q{计算：$\displaystyle \frac{5}{11} \times 79\frac{18}{19} + 40 \times \frac{1}{11} + \frac{1}{11} \times \frac{5}{19} + 5\frac{1}{4} \div \frac{7}{24}$}

\Q{（工程问题）一项工程，甲工程队做6天休息1天，乙工程队做4天休息2天。甲工程队单独做完要32天，乙工程队单独做完要34天。现在两队合作，那么完成这项工程要\underline{\hspace{3em}}天。}

\Qlast{（容斥原理）某校 250 名学生在一次语、数、外三科竞赛中，参加语文竞赛的有 120 人，参加数学竞赛的有 130 人，参加外语竞赛的有 100 人，既参加语文又参加数学竞赛的有 45 人，既参加数学又参加外语竞赛的有 50 人，既参加语文又参加外语竞赛的有 40 人，有 15 人这三项竞赛都不参加。则三项都参加的共有多少人？}

\end{OnePageFourFill}

% =========================================================
% 7.3
% =========================================================
\section*{7月3日作业（4道题当晚22:00前打卡）}
\begin{OnePageFourFill}

\Q{计算：$\displaystyle 2025 \div 2025\frac{2025}{2026} \div \frac{2026 + \dfrac{1}{2027}}{2027 + \dfrac{1}{2026}}$}

\Q{计算：$\displaystyle \frac{4}{7} \times 69\frac{14}{15} + 70 \times \frac{1}{7} + \frac{1}{7} \times \frac{4}{15} + 2\frac{4}{9} \div \frac{11}{45}$}

\Q{（工程问题）一项工程，甲工程队做5天休息2天，乙工程队做2天休息1天。甲工程队单独做完要20天，乙工程队单独做完要23天。现在两队合作，那么完成这项工程要\underline{\hspace{3em}}天。}

\Qlast{（容斥原理）某校 300 名学生在一次语、数、外三科竞赛中，参加语文竞赛的有 140 人，参加数学竞赛的有 150 人，参加外语竞赛的有 130 人，既参加语文又参加数学竞赛的有 55 人，既参加数学又参加外语竞赛的有 60 人，既参加语文又参加外语竞赛的有 50 人，有 20 人这三项竞赛都不参加。则三项都参加的共有多少人？}

\end{OnePageFourFill}

\end{document}
